\subsection*{Estudio legal}
\addcontentsline{toc}{subsection}{Estudio legal}

El estudio legal de este proyecto, y por ende de la empresa propuesta, tiene como objetivo analizar la viabilidad del negocio frente a las normas y leyes establecidas dentro del marco jurídico colombiano. En este apartado se identifican los aspectos legales relacionados con la creación de empresas, sus alcances, limitaciones y obligaciones, de acuerdo con la naturaleza del proyecto. Asimismo, se busca describir de manera clara el proceso de formalización de la empresa ante las entidades estatales correspondientes, garantizando que su funcionamiento se desarrolle dentro del marco legal vigente y cumpla con las regulaciones requeridas.

\subsubsection*{Tipo de sociedad}
\addcontentsline{toc}{subsubsection}{Tipo de sociedad}

Nuestra empresa será constituida bajo la figura de Sociedad por Acciones Simplificada (S.A.S.), tipo societario creado en la legislación colombiana mediante la Ley 1258 de 2008. Este modelo se ha convertido en uno de los más utilizados en el país debido a las ventajas que ofrece en términos de simplicidad administrativa y flexibilidad organizacional. Adicionalmente, para efectos tributarios, las S.A.S. se rigen por las normas aplicables a las sociedades anónimas \parencite{Ley1258}.

Este tipo de sociedad es el mejor para nuestra empresa, especialmente debido a las ventajas que trae dentro del ámbito tributarios. Ya que desde la entrada en vigencia de la Ley 1258, alrededor del 54\% de las empresas colombianas lo hacen bajo esta modalidad \parencite{Actualicese01}. La elección de esta sociedad nos permite simplificar los trámites burocráticos, de esta forma disminuyendo los costos del inicio y facilitando el proceso de formalización de nuestra empresa. Junto a ello, las S.A.S. no necesitan un revisor fiscal, permitiendo un ahorro significativo, y estas sociedades pueden ser llevadas a cabo tanto por personas naturales como jurídicas, brindando flexibilidad hacia los emprendedores.

\subsubsection*{Especificaciones generales}
\addcontentsline{toc}{subsubsection}{Especificaciones generales}

Nuestra empresa será constituida bajo las siguientes características generales:

\begin{itemize}
	\item \textbf{Socios:} Leidy Marcela Morales Segura con Cédula de Ciudadanía No. y David Leonardo Rincón Ovalle con Cédula de Ciudadanía No. 1001093774.
	
	\item \textbf{Nombre de la empresa:} EduFinanzas.
	
	\item \textbf{Duración:} Tiempo indefinido.
	
	\item \textbf{Objeto social:} Desarrollo e innovación de herramientas tecnológicas orientadas a la educación financiera.
	
	\item \textbf{Responsabilidad sobre los aportes:} 50\% por cada socio.
	
	\item \textbf{Representante legal:} David Leonardo Rincón Ovalle.  
\end{itemize}

\subsubsection*{Obligaciones legales}
\addcontentsline{toc}{subsubsection}{Obligaciones legales}

Las obligaciones legales descritas a continuación se encuentran enmarcadas principalmente en la Ley 1258 de 2008, mediante la cual se regula la Sociedad por Acciones Simplificada (S.A.S.) \parencite{Ley1258}. Esta normativa establece las responsabilidades y deberes que debe cumplir la empresa para operar legalmente bajo esta figura jurídica en Colombia.

\begin{enumerate}
	\item Realizar la Asamblea Ordinaria de Accionistas al menos una vez al año.

	\item Renovar la matrícula mercantil dentro de los primeros tres (3) meses de cada año ante la Cámara de Comercio.

	\item Disponer de firma electrónica para los representantes legales, de acuerdo con la Resolución 070 de la Dirección de Impuestos y Aduanas Nacionales (DIAN) \parencite{Resolucion70}.

	\item Cumplir con las obligaciones como agentes de retención en la fuente a título de renta, IVA, ICA, entre otros impuestos aplicables.

	\item Expedir facturación conforme a las normas establecidas por la DIAN.

	\item Cumplir con el Gravamen a los Movimientos Financieros (GMF) cuando corresponda.

	\item Llevar la contabilidad de la empresa conforme a las normas contables vigentes.

	\item Contar con revisor fiscal cuando los ingresos o activos de la empresa superen los montos establecidos por la ley.
\end{enumerate}

\subsubsection*{Proceso para disolución de la empresa}
\addcontentsline{toc}{subsubsection}{Proceso para disolución de la empresa}

En caso de disolución de la empresa constituida bajo la figura de Sociedad por Acciones Simplificada (S.A.S.), se deberán seguir los procedimientos establecidos por la normativa legal colombiana, los cuales incluyen las siguientes etapas:

\begin{enumerate}
	\item En el casos de las S.A.S., Empresa Unipersonal y Sociedades de la Ley 1014 de 2006, la declaratoria de disolución por mutuo acuerdo podrá realizarse mediante acta o documento privado firmado por las partes involucradas y aprobado por la Asamblea de Accionistas. Posteriormente este documento será inscrito en el registro mercantil de la Cámara de Comercio \parencite{Ley1014}.
	
	\item Posteriormente, será necesario realizar registrar el acta de disolución ante la Cámara de Comercio. A partir de este momento, la empresa pasará a tener el sufijo ``en liquidación''. Se entregará copia del acta en la Cámara de Comercio, mientras que el acta original quedará en los archivos de nuestra empresa.
	
	\item El siguiente paso es informar a la Dirección de Impuestos y Aduanas Nacionales (DIAN) sobre la existencia de posibles deudas fiscales. Este reporte debe realizarse dentro de los diez (10) días posteriores al registro de la disolución.
	
	\item Acto seguido se publican avisos que informen a terceros sobre el inicio del proceso de liquidación de la sociedad.
	
	\item Se elabora un inventario del patrimonio social y el balance final de la empresa, detallando todos los activos y pasivos.
	
	\item Luego se paga el pasivo externo, es decir, el pasivo externo, es decir, las deudas con terceros, incluyendo obligaciones fiscales pendientes y la declaración de renta final.
	
	\item Una vez pagadas las deudas, el liquidador procederá a distribuir los remanentes entre los socios o accionistas de la sociedad.
	
	\item Después se elabora el proyecto de liquidación, el cual deberá incluir:
	
	\begin{itemize}
		\item Inventarios.
		
		\item Balance general.
		
		\item Estado de pérdidas y ganancias.
		
		\item Relación de los pasivos de la entidad.
		
		\item Estado de cuentas a la DIAN.
		
		\item Pago de pasivos.
		
		\item Indicación del remanente.
		
		\item Destinación del remanente.
	\end{itemize}
	
	\item Luego se convoca una reunión e la Asamblea de Accionistas para aprobar el proyecto de liquidación presentado por el liquidador.
	
	\item Finalmente, se hace registro del acta de la cuenta final de liquidación ante la Cámara de Comercio. Se deberá pagar un impuesto de registro equivalente al 0.7\% sobre el valor de los remanentes de la empresa, después de pagar el pasivo externo. En caso de no existir remanentes, el trámite se realizará como acto sin cuantía.
\end{enumerate}

\subsubsection*{Propiedad intelectual}
\addcontentsline{toc}{subsubsection}{Propiedad intelectual}

De acuerdo al Decreto 1360 de 1989, artículo 1, presenta lo siguiente: ``De conformidad con lo previsto en la Ley 23 de 1982 sobre Derechos de Autos, el soporte lógico (software) se considera como creación propia del dominio literario'' \parencite{Ley1360}. En consecuencia, el software desarrollado en el marco de este proyecto deberá ser registrado como obra ante la Dirección Nacional de Derechos de Autor, con el fin de proteger la propiedad intelectual del producto tecnológico desarrollado.

\subsubsection*{Permisos y licencias requeridos}
\addcontentsline{toc}{subsubsection}{Permisos y licencias requeridos}

Para la operación legal del negocio, se requiere cumplir con diversos registros y permisos ante las autoridades correspondientes. Aunque la naturaleza del proyecto no corresponde exactamente a un comercio electrónico tradicional, al ofrecer servicios digitales en línea comparte algunas obligaciones regulatorias asociadas a este tipo de actividades.

\begin{itemize}
	\item \textbf{Registro de Unidad Tributaria (RUT):} Registro obligatorio ante la DIAN para obtener una Unidad Tributaria, que es el número de identificación fiscal de nuestra compañia.
	
	\item \textbf{Registro de Empresas Unificadas de Servicios (RUES):} Registro obligatorio administrado por la Cámara de Comercio para poder inscribir la empresa en la sociedad del Registro Mercantil.
	
	\item \textbf{Registro del Sistema de Administración de Riesgo de Lavado de Activos y Financiación del Terrorismo (SARLAFT):} Registro realizado ante la Superintendencia Financiera bajo el cual se específica que buscamos prevenir que las actividades de nuestra empresas no se utilicen para lavar activos o financiar terroristas. Este registro es obligatorio para cualquier empresa (que posiblemente presenten características financieras) que busque operar en Colombia.
	
	\item \textbf{Registro de Comercio:} Es el registro que se hace en la Cámara de Comercio con el objetivo de tener una licencia de comercio. Esta licencia es obligatoria para cualquier empresa que desee operar en Colombia.
	
	\item \textbf{Registro Mercantil:} Inscripción de la empresa ante la Cámara de Comercio que acredita legalmente su existencia y actividad comercial, la cual es sobre la venta de servicios o suscripciones.
	
	\item \textbf{Sistema de Gestión de Seguridad y Salud en el Trabajo (SG-SST):} Registro ante el Ministerio de Trabajo para garantizar condiciones adecuadas de seguridad y salud laboral.
	
	\item \textbf{Registro de bases de datos personales:} Registro ante la Superintendencia de Industria y Comercio para cumplir con la normativa de protección de datos personales.
	
	\item \textbf{Cumplimiento de normativa de comercio electrónico:} Regulación supervisada por la Superintendencia de Industria y Comercio para empresas que prestan servicios digitales en línea.
\end{itemize}

\subsubsection*{Normas técnicas y de calidad}
\addcontentsline{toc}{subsubsection}{Normas técnicas y de calidad}

Para garantizar la calidad del servicio, la seguridad de la información y la confiabilidad del sistema, la empresa buscará alinearse con diversas normas técnicas y estándares internacionales:

\begin{itemize}
	\item \textbf{Norma Técnica Colombiana (NTC) ISO 9000:} Describe los fundamentos sobre los sistemas de gestión de la calidad y especifica la terminología de los sistemas de gestión de la calidad. Esta norma es fundamental para garantizar la calidad y satisfacción del cliente.
	
	\item \textbf{Norma ISO 9001:} Es el estándar de calidad más aplicado dentro de todo el mundo. Acá se describen la implementación de los procesos y procedimientos destinados para lograr la satisfacción del cliente. La certificación ISO 9001 demuestra que EduFinanzas tiene un compromiso para garantizar la calidad, la consistencia y la autenticidad.
	
	\item \textbf{Norma ISO 10008:} Es el estándar de gestión específico para el comercio electrónico. Asegura que se implementen los procesos y tareas con el objetivo de proteger la información y datos de nuestros clientes.
	
	\item \textbf{Norma ISO 27001:} Es un estándar de seguridad de la información que busca garantizar la protección de la información y los datos de nuestro clientes. Esta norma es esencial para proteger la información y datos de EduFinanzas.
	
	\item \textbf{Norma ISO 20000:} Es un estándar de gestión de continuidad de los negocios que busca garantizar la continuidad de los negocios en caso de que ocurran desastres y situaciones críticas. Esta norma es esencial para proteger la continuidad de los negocios.
	
	\item \textbf{Norma ISO 22301:} Es un estándar de gestión para continuidad de los negocios que busca garantizar la continuidad de los negocios en caso de que ocurran desastres o situaciones críticas. Esta norma es fundamental para proteger la continuidad de los negocios.
	
	\item \textbf{Norma ISO 27002:} Es el estándar de seguridad de la información que proporciona recomendaciones prácticas para implementar medidas de seguridad en la información de un negocio.
\end{itemize}

